\section[Introdução]{Introdução}

O Projeto foi concebido através de uma ``dor'' que o stakeholder, Lucas Savaris, e pessoas próximas a ele tinham: o mau aproveitamento de seu tempo por conta das viciantes redes sociais.

O projeto tem como objetivo auxiliar o usuário a reduzir o tempo de uso do celular e tornar o seu uso mais produtivo. Com esse objetivo, foi proposto um aplicativo móvel que permite ao usuário definir o tempo que deseja passar utilizando o celular, e receber notificações personalizadas quando atingir o limite estabelecido.

Além disso, o aplicativo também apresentará uma barra de progresso indicando o tempo restante com base no tempo de uso definido pelo usuário. Caso o usuário ultrapasse o limite estabelecido, o aplicativo ativará uma terceira funcionalidade denominada \textit{Greyscale}, na qual a tela do celular ficará em escala de cinza.

A execução do projeto ocorreu no primeiro semestre de 2023, entre os dias 27 de fevereiro e 21 de junho, pelos estudantes do bacharelado de Engenharia de Software em diversas etapas da \ac{ages}, com orientação do professor Marcelo Yamaguti. A foto do time responsável pelo projeto pode ser vista na Figura~\ref{fig:focus-time}.

\begin{figure}[H]
    \centering
    \small
    \caption{Time do projeto Focus}
    \fbox{\parbox[c][7cm][c]{0.9\linewidth}{\centering \textit{[Inserir imagem]}}}

    Fonte: Wiki do projeto \cite{wikiFocus}
    \label{fig:focus-time}
\end{figure}